% Compile with: uplatex file.tex && dvipdfmx file.dvi
%           or: latexmk -uplatex file.tex
%\documentclass[uplatex,dvipdfmx]{jsarticle}
\documentclass[dvipdfmx,a4paper,11pt]{jsarticle}
\usepackage{amsmath, amssymb}
\usepackage{tikz}

\title{scenario\_rlc\_transfer\_function}
\date{\today}

\begin{document}
\maketitle

直列RLC回路（電荷 $q(t)$、入力電圧 $v(t)$）を扱う。$q,v$ を時間依存関数として準備する。

q を時間依存関数として定義する

$\checkmark$ q を時間依存として宣言しました（時間変数: t）

v を時間依存関数として定義する

$\checkmark$ v を時間依存として宣言しました（時間変数: t）

$q$ のラプラス変換を $Q$（s領域の量）と名付ける定義カードを作る。

q のラプラス変換を Q とおく

\begin{equation}
  Q(s) = \int_0^{\infty} q(t)\, e^{-s t}\, dt
  \label{eq:E1}
\end{equation}

同様に $v$ のラプラス変換を $V$ と名付ける。

v のラプラス変換を V とおく

\begin{equation}
  V(s) = \int_0^{\infty} v(t)\, e^{-s t}\, dt
  \label{eq:E2}
\end{equation}

導関数のラプラス変換則を書く。部分積分は $u=e^{-st}$ が変換の構造で決まっていて選択が無いので、これは計算ではなく定理であり、研究者が書ける。まず $\mathcal L\{\dot q\} = sQ - q(0)$。

$s Q - q_0$

\begin{equation}
  s Q - q_0
  \label{eq:E3}
\end{equation}

続いて $\mathcal L\{\ddot q\} = s^2 Q - s q(0) - \dot q(0)$。

$s^2 Q - s q_0 - u_0$

\begin{equation}
  s^2 Q - s q_0 - u_0
  \label{eq:E4}
\end{equation}

回路方程式 $L\ddot q+R\dot q+q/C=v$（インダクタンス $L$・抵抗 $R$・キャパシタンス $C$）の $s$ 領域形を、上の2つを E 番号で参照して組み立てる。

$L E4 + R E3 + Q/C = V$

\begin{equation}
  L \left(-u_{0}+Q s^2-q_{0} s\right)+R \left(Q s-q_{0}\right)+\frac{Q}{C}=V
  \label{eq:E5}
\end{equation}

伝達関数はゼロ初期条件のもとで定義される量なので、$q(0)=0$、$\dot q(0)=0$ を置く。H-1 の定義は登録時に CAS へ送られるので、これ以降の計算に自動的に効く。

$q_0 = 0$

\begin{equation}
  q_0 = 0
  \label{eq:E6}
\end{equation}

$u_0 = 0$

\begin{equation}
  u_0 = 0
  \label{eq:E7}
\end{equation}

伝達関数 $H=Q/V$（出力と入力の $s$ 領域での比）を定義する。

$H = Q/V$

\begin{equation}
  H = Q/V
  \label{eq:E8}
\end{equation}

$s$ 領域の方程式を $Q$ について解く。ゼロ初期条件は既に効いている。

E5 を Q について解いて

\begin{equation}
  Q = \frac{C V}{C L s^2+C R s+1}
  \label{eq:E9}
\end{equation}

得られた $Q$ の式を $H$ の定義に代入し、伝達関数の閉じた式を得る。

E9 を E8 に代入して

\begin{equation}
  H = \frac{C}{C L s^2+C R s+1}
  \label{eq:E10}
\end{equation}

\end{document}